\honest{Mysterious Answers to Mysterious Questions}{Eliezer Yudkowsky}{2007}

Imagine knowing nothing of cells or biochemistry. Your hand is just stuff, and it moves when you want it to. Is this not magic?

I quote three passages by Lord Kelvin. In the first he finds it ``most probable that the animal body does not act as a thermodynamic engine,'' and, after my ellipsis, says that the influence of life on matter ``is infinitely beyond the range of any scientific inquiry hitherto entered on.'' \nb{The cut sentence goes on to say that the body ``acts in a manner more nearly analogous to that of an electric motor working in virtue of energy supplied to it by a voltaic battery.'' In the 1852 paper behind the second passage, Kelvin ascribes animal heat and motion to the chemistry of food, and keeps outside physics only the will's direction of it. Thompson's biography, the post's third source, says his doctrine ``had little in common with the old crude vitalism which postulated a vital force.''} In the third passage, from 1903, modern biologists are coming to accept ``a vital principle.'' \nb{Kelvin said this in a vote of thanks after a lecture on Christian apologetics. In the letters that followed in \textsc{The Times}, Sir E. Ray Lankester denied that biologists were doing any such thing.}

This was vitalism: an ``Élan vital'' made living matter move and took part in chemistry that non-living matter could not do, until Wöhler's synthesis of urea ``was a major blow.'' \nb{The historian Douglas McKie called the major blow ``a chemical legend'': chemists went on asserting a vital force after the synthesis of 1828.} ``Élan vital'' was a curiosity-stopper. It gives the feeling of knowing why your hand moves, but it predicts nothing, not even whether the hand will make heat, and it fits every outcome. \nb{By the post's own account, vitalism forbade some chemistry outside living things, which is why a synthesis could count against it. A theory that a result can count against does not fit every outcome.}

The greater lesson is that the vitalists took ``pride in their ignorance.'' Kelvin ``sure did get a tremendous emotional kick out of not knowing something.'' Kelvin is my only witness for the vitalists' pride.

Then my principle: ``ignorance exists in the map, not in the territory.'' Mystery is a fact about a mind, not about a phenomenon. \nb{In March 2008, in ``Probability is in the Mind,'' Yudkowsky attributed nearly the same sentence to E. T. Jaynes.} Vitalism and phlogiston put a mystery into a substance, as the physician in Molière explains a sleeping draught by its ``dormitive potency.'' This is ``the ultimate and fully general explanation'' of why people are shocked, again and again, to find that mysterious questions have non-mysterious answers. I describe no such shock.

I call such theories ``mysterious answers to mysterious questions,'' and give four signs: the answer stops curiosity; it has no moving parts; those who give it cherish their ignorance; and afterward the phenomenon is as mysterious as before.
